\documentclass[12pt]{article}
\usepackage[utf8]{inputenc}
\usepackage[T1]{fontenc}
\usepackage{hyperref}
\title{Kaki Ae Language Resources for K-12 Curriculum Development}
\author{Compiled by Claude Sonnet 5 (claude-sonnet-5) for Andr\'as Kornai}
\date{}
\begin{document}
\maketitle

\section{Introduction}
Kaki Ae (ISO 639-3 tbd; also called Tate) is spoken near Kerema in Gulf Province, Papua New Guinea. Sources give about 300 to 630 speakers. Its classification is disputed: it may be related to the Eleman languages, or it may be an isolate. Searches found only linguistic and reference material. They found no school materials, primers or textbooks.

\section{Primary Educational Resources}
No primers, textbooks or school programmes for Kaki Ae turned up in search results. A curriculum would have to be built from the linguistic documentation below.

\section{Linguistic and Reference Materials}
\begin{itemize}
\item John M. Clifton, ``A grammar sketch of the Kaki Ae language'', SIL Work Papers vol.\ 39. The search summary says it covers phonology, phrases and morphemes, and includes a Kaki Ae--English lexicon, an English--Kaki Ae finder list and seven short texts. \url{https://commons.und.edu/sil-work-papers/vol39/iss1/2}
\item Clifton, John M.\ (1997), \emph{The Kaki Ae language}, Pacific Linguistics, listed in the Glottolog reference record. \url{https://glottolog.org/resource/reference/id/79499}
\item Wikipedia article on Kaki Ae. It has a village list, pronouns and phonology. \url{https://en.wikipedia.org/wiki/Kaki_Ae_language}
\item ASJP word-list entry for Kaki Ae. \url{https://asjp.clld.org/languages/KAKI_AE}
\item Numerals data page. \url{https://lingweb.eva.mpg.de/channumerals/Kaki-Ae.htm}
\end{itemize}

\section{Religious and Community Resources}
None were found in search results.

\section{Digital and Audio Resources}
No audio, video or software resources were found.

\section{Curriculum Development Context}
Little exists for this language beyond Clifton's descriptive work. Any K-12 curriculum would need new materials developed with the community. Clifton's lexicon and texts are the most likely starting point for an orthography and vocabulary.
\end{document}